\documentclass[11pt]{article}
\usepackage[margin=1in]{geometry}
\usepackage{amsmath,amssymb,amsthm,mathtools}
\usepackage{microtype}
\usepackage[hidelinks]{hyperref}
\usepackage{enumitem}
\usepackage{booktabs}
\usepackage{bm}
\usepackage{authblk}

\newtheorem{theorem}{Theorem}[section]
\newtheorem{proposition}[theorem]{Proposition}
\newtheorem{corollary}[theorem]{Corollary}
\newtheorem{lemma}[theorem]{Lemma}
\theoremstyle{remark}
\newtheorem{remark}[theorem]{Remark}

\newcommand{\Hsq}{H^2}
\newcommand{\Ran}{\operatorname{Ran}}
\newcommand{\spec}{\sigma}
\newcommand{\cF}{c_F}
\newcommand{\1}{\mathbf 1}

\title{Principal Angles Between Arithmetic Weighted-Composition Ranges in $H^2$}
\author{Shad Claiborne\\\small Independent Researcher\\\small ORCID: \href{https://orcid.org/0009-0005-9443-3679}{0009-0005-9443-3679}}
\date{Preprint v1.0\\\small DOI: \href{https://doi.org/10.5281/zenodo.22885535}{10.5281/zenodo.22885535}}

\begin{document}
\maketitle

\begin{abstract}
Let
\[
W_nf(z)=(1+z+\cdots+z^{n-1})f(z^n),\qquad V_n=n^{-1/2}W_n,
\]
on the Hardy space $H^2$ of the unit disk, and set $\mathcal H_n=\Ran V_n$.
For coprime integers $2\le m<n$, we reduce the compression geometry of
$\mathcal H_m$ and $\mathcal H_n$ to a finite weighted path Laplacian whose
edge weights are determined by the residues $jn\pmod m$. This yields exact
sum and product identities for the principal-angle defects. We prove that the
algebraic connectivity of the resulting path Laplacian is at most $2$, with
equality exactly when $n\equiv\pm1\pmod m$. In those extremal cases the
Laplacian has the classical Hahn spectrum $k(k+1)$, and hence the squared
cosines of the distinct principal-angle values are
$1-k(k+1)/(mn)$, $0\le k<m$. We also record the corresponding Friedrichs-angle
and alternating-projection consequences. All results are unconditional.
\end{abstract}

\section{Introduction}
For $n\ge1$, consider the weighted composition operators
\[
W_nf(z)=(1+z+\cdots+z^{n-1})f(z^n)
\]
on $\Hsq$. The semigroup $(W_n)_{n\ge1}$ and its connections with the
Riemann hypothesis, invariant-subspace questions, and periodic dilation
completeness were studied by Manzur Villa, Noor, and Santos
\cite{ManzurNoorSantos2023}. In particular, $W_n^*W_n=nI$, so
$V_n=W_n/\sqrt n$ is an isometry.

This note studies a different feature of the same semigroup: the mutual
positions of the closed ranges
\[
\mathcal H_n:=\Ran V_n.
\]
Principal angles and the related Friedrichs angle provide standard measures
of the geometry of two closed subspaces; see, for example,
Bjorck--Golub \cite{BjorckGolub1973}. A recent RH-related use of Friedrichs
angles, for different subspaces in the Nyman--Beurling setting, appears in
Yang \cite{Yang2026}.

Our main observation is that for coprime $m,n$ the infinite-dimensional
problem decomposes into copies of one finite $mn$-coordinate problem. The
finite cross-Gram matrix is governed by a weighted path Laplacian. The
weights are a permutation of
\[
\{j(m-j):1\le j<m\},
\]
with their order determined arithmetically by multiplication by $n$ modulo
$m$. This gives two order-independent spectral invariants and a sharp
extremal theorem. In the congruence classes $n\equiv\pm1\pmod m$ the path is
the classical symmetric Hahn difference operator, giving a complete closed
form for the spectrum.

\section{Coefficient model and block-constant ranges}
Identify $\Hsq$ isometrically with $\ell^2(\mathbb N_0)$ via Taylor
coefficients, and let $(e_j)_{j\ge0}$ denote its standard basis. Then
\[
V_ne_j=\frac1{\sqrt n}\sum_{r=0}^{n-1}e_{nj+r}.
\]
Thus $\mathcal H_n$ consists exactly of sequences that are constant on each
consecutive block of $n$ coordinates. Let $P_n=V_nV_n^*$ be the orthogonal
projection onto $\mathcal H_n$.

\begin{proposition}\label{prop:intersection}
If $\gcd(m,n)=1$, then
\[
\mathcal H_m\cap\mathcal H_n=\mathcal H_{mn}.
\]
\end{proposition}

\begin{proof}
Work inside a block of $mn$ consecutive coordinates. A vector in the
intersection is constant on every $m$-block and on every $n$-block. The
bipartite overlap graph of these two partitions is connected: an interior
cut separating the overlap graph would have to occur at a coordinate that is
simultaneously a multiple of $m$ and of $n$, and coprimality excludes such a
coordinate before $mn$. Hence the vector is constant on the entire
$mn$-block. Different $mn$-blocks are independent, proving the claim.
\end{proof}

The same argument shows that the pair $(\mathcal H_m,\mathcal H_n)$ is an
orthogonal infinite repetition of a single finite pair of block-constant
subspaces in $\mathbb C^{mn}$.

\section{The arithmetic overlap Laplacian}
Assume throughout this section that $2\le m<n$ and $\gcd(m,n)=1$. Inside
one $mn$-block, let $u_j$ ($0\le j<m$) be the normalized indicators of the
$n$-blocks and $v_i$ ($0\le i<n$) the normalized indicators of the
$m$-blocks. Let
\[
C_{ij}=\langle v_i,u_j\rangle.
\]
The eigenvalues of $G=C^*C$ are the squared singular values of the
cross-projection and hence the squared cosines of the finite-block principal
angles \cite{BjorckGolub1973}.

For $1\le j<m$, let $r_j$ be the least positive residue of $jn$ modulo $m$:
\[
r_j\equiv jn\pmod m,\qquad 1\le r_j\le m-1,
\]
and put
\[
w_j=r_j(m-r_j).
\]
Let $L_{m,n}$ be the weighted path Laplacian on $m$ vertices with edge
weights $w_1,\ldots,w_{m-1}$, i.e.
\[
(L_{m,n})_{j,j+1}=(L_{m,n})_{j+1,j}=-w_{j+1}
\]
with diagonal entries equal to the sum of incident edge weights.

\begin{theorem}[Coprime compression reduction]\label{thm:compression}
For coprime $2\le m<n$,
\[
G=I-\frac1{mn}L_{m,n}.
\]
Consequently, if
\[
0=\mu_0<\mu_1<\cdots<\mu_{m-1}
\]
are the eigenvalues of $L_{m,n}$, then the distinct blockwise
principal-angle values satisfy
\[
\cos^2\theta_k=1-\frac{\mu_k}{mn},\qquad 0\le k<m.
\]
Each value occurs with infinite multiplicity for the full pair
$(\mathcal H_m,\mathcal H_n)$.
\end{theorem}

\begin{proof}
The boundary between consecutive $n$-blocks occurs at $jn$. Relative to the
$m$-partition, exactly one $m$-block straddles that boundary. Its two pieces
have lengths $r_j$ and $m-r_j$, so the off-diagonal contribution to the
neighboring Gram entry is
\[
\frac{r_j(m-r_j)}{mn}=\frac{w_j}{mn}.
\]
Because $n>m$, no single $m$-block meets two nonadjacent $n$-blocks. Hence
$G$ is tridiagonal, with neighboring entries $w_j/(mn)$ and diagonal entries
$1$ minus the sum of the adjacent normalized edge weights. This is exactly
$I-L_{m,n}/(mn)$.

All $w_j$ are positive, so the weighted path is connected. Its Laplacian is
an irreducible symmetric tridiagonal matrix and therefore has simple
spectrum. The principal-angle statement follows from the standard
cross-Gram characterization. Finally, the same $mn$-block geometry repeats
orthogonally along $\ell^2(\mathbb N_0)$, giving infinite multiplicity.
\end{proof}

\section{Universal spectral invariants}
Since multiplication by $n$ permutes the nonzero residues modulo $m$,
\[
\{r_1,\ldots,r_{m-1}\}=\{1,\ldots,m-1\}.
\]
Thus the multiset of edge weights is independent of the residue class of
$n$; only their order changes.

\begin{theorem}[Trace and product invariants]\label{thm:invariants}
Under the hypotheses of Theorem~\ref{thm:compression},
\[
\sum_{k=1}^{m-1}\sin^2\theta_k=\frac{m^2-1}{3n},
\]
and
\[
\prod_{k=1}^{m-1}\sin^2\theta_k
=\frac{m((m-1)!)^2}{(mn)^{m-1}}.
\]
\end{theorem}

\begin{proof}
First,
\[
\sum_{j=1}^{m-1}w_j
=\sum_{r=1}^{m-1}r(m-r)
=\frac{m(m^2-1)}6.
\]
The trace of a weighted path Laplacian is twice the total edge weight, so
\[
\sum_{k=1}^{m-1}\mu_k=\operatorname{tr}L_{m,n}
=\frac{m(m^2-1)}3.
\]
Using $\sin^2\theta_k=\mu_k/(mn)$ gives the sum identity.

For the product, the weighted matrix-tree theorem applied to a path gives
\[
\prod_{k=1}^{m-1}\mu_k=m\prod_{j=1}^{m-1}w_j.
\]
Since the $r_j$ form a permutation of $1,\ldots,m-1$,
\[
\prod_{j=1}^{m-1}w_j
=\prod_{r=1}^{m-1}r(m-r)=((m-1)!)^2.
\]
Dividing by $(mn)^{m-1}$ yields the second formula.
\end{proof}

\section{A sharp extremal theorem}
We next determine the largest possible algebraic connectivity among the
arithmetic orderings above.

\begin{theorem}[Sharp algebraic-connectivity bound]\label{thm:extremal}
Let $0=\mu_0<\mu_1<\cdots<\mu_{m-1}$ be the spectrum of $L_{m,n}$. Then
\[
\mu_1\le2.
\]
Equality holds if and only if
\[
n\equiv\pm1\pmod m.
\]
\end{theorem}

\begin{proof}
Let
\[
x_j=j-\frac{m-1}{2},\qquad 0\le j<m.
\]
Then $x\perp\1$, and $x_j-x_{j-1}=1$. The Rayleigh quotient is therefore
\[
\frac{x^TL_{m,n}x}{x^Tx}
=\frac{\sum_{j=1}^{m-1}w_j}
{\sum_{j=0}^{m-1}(j-(m-1)/2)^2}=2,
\]
using the sum computed in Theorem~\ref{thm:invariants}. By Rayleigh--Ritz,
$\mu_1\le2$.

If $\mu_1=2$, then $x$ attains the minimum and hence
$L_{m,n}x=2x$. At the left endpoint,
\[
-w_1=2x_0=-(m-1),
\]
so $w_1=m-1$. Since $w_1=r_1(m-r_1)$, this implies
$r_1=1$ or $r_1=m-1$. But $r_1\equiv n\pmod m$, hence
$n\equiv\pm1\pmod m$.

Conversely, if $n\equiv\pm1\pmod m$, then
$r_j=j$ or $m-j$, so in either case $w_j=j(m-j)$. The Hahn-spectrum
calculation in Section~\ref{sec:hahn} gives $\mu_1=2$.
\end{proof}

\begin{remark}
The equality argument also determines the entire positional weight sequence.
Indeed, if $Lx=2x$ for the centered linear vector, then the interior vertex
equations give
\[
w_{j+1}=w_j+(m-1)-2j,
\]
and induction from $w_1=m-1$ yields $w_j=j(m-j)$.
Thus among all permutations of the fixed multiset
$\{j(m-j):1\le j<m\}$, the natural Hahn ordering is the unique positional
weight sequence attaining algebraic connectivity $2$ (its reversal is the
same sequence because $j(m-j)=(m-j)j$).
\end{remark}

\begin{corollary}[Friedrichs-angle extremality]\label{cor:friedrichs}
For coprime $2\le m<n$,
\[
\cF(\mathcal H_m,\mathcal H_n)
\ge \sqrt{1-\frac{2}{mn}},
\]
with equality if and only if $n\equiv\pm1\pmod m$.
\end{corollary}

\begin{proof}
After removing the common intersection, the largest squared cosine is
$1-\mu_1/(mn)$. Apply Theorem~\ref{thm:extremal}.
\end{proof}

\section{The Hahn closed-form cases}\label{sec:hahn}
If $n\equiv\pm1\pmod m$, then
\[
w_j=j(m-j),\qquad 1\le j<m.
\]
This weighted second-order difference operator is the symmetric Hahn
operator. Its classical spectrum is
\[
\{k(k+1):0\le k<m\};
\]
see, for example, the Hahn-polynomial treatment in
\cite{KoekoekLeskySwarttouw2010}.

For completeness, there is also a short elementary verification. Identify a
vector in $\mathbb C^m$ with the values of a polynomial $p(j)$ of degree at
most $m-1$. Then
\[
(L_mp)(j)=j(m-j)[p(j)-p(j-1)]
 +(j+1)(m-j-1)[p(j)-p(j+1)].
\]
If $p(j)=aj^k$ plus lower-order terms, the degree-$k+1$ terms cancel and the
leading surviving term is $ak(k+1)j^k$. Hence the operator preserves the
polynomial-degree filtration and is triangular in the monomial basis with
diagonal $k(k+1)$.

\begin{theorem}[Complete extremal spectrum]\label{thm:hahn}
If $2\le m<n$ and $n\equiv\pm1\pmod m$, then
\[
\spec(P_nP_mP_n|_{\mathcal H_n})
=\left\{1-\frac{k(k+1)}{mn}:0\le k<m\right\}.
\]
Equivalently, the distinct principal-angle values satisfy
\[
\cos^2\theta_k=1-\frac{k(k+1)}{mn},\qquad 0\le k<m.
\]
Each value occurs with infinite multiplicity globally.
\end{theorem}

\section{Alternating projections}
For two closed subspaces $M,N$, the exact two-subspace alternating-projection
formula of Kayalar and Weinert \cite{KayalarWeinert1988} is
\[
\|(P_NP_M)^r-P_{M\cap N}\|=\cF(M,N)^{2r-1},\qquad r\ge1.
\]
Combining this with Proposition~\ref{prop:intersection} and
Corollary~\ref{cor:friedrichs} gives the following.

\begin{corollary}
For coprime $2\le m<n$,
\[
\|(P_nP_m)^r-P_{mn}\|
=\left(1-\frac{\mu_1}{mn}\right)^{r-1/2}.
\]
In the extremal cases $n\equiv\pm1\pmod m$,
\[
\|(P_nP_m)^r-P_{mn}\|
=\left(1-\frac2{mn}\right)^{r-1/2}.
\]
\end{corollary}

\section{Discussion}
The coprime reduction separates two pieces of arithmetic information. The
multiset of edge weights is universal for fixed $m$, while their order is
determined by the residue class of $n$ modulo $m$. The trace and
pseudodeterminant are therefore order-independent, whereas the individual
principal angles generally depend on the arithmetic ordering. The
congruence classes $n\equiv\pm1\pmod m$ are distinguished twice: they
maximize the relevant spectral separation and simultaneously admit the
complete Hahn closed form.

The weighted-composition semigroup was developed in RH-related work
\cite{ManzurNoorSantos2023}, but no statement in this paper assumes the
Riemann hypothesis or asserts a proof of it.

\section*{Research assistance note}
ChatGPT (OpenAI) was used for editorial assistance, literature-search support,
and numerical verification. The author is solely responsible for all
mathematical statements, proofs, and conclusions.

\begin{thebibliography}{9}
\bibitem{ManzurNoorSantos2023}
J. C. Manzur Villa, S. W. Noor, and C. F. dos Santos,
\emph{A weighted composition semigroup related to three open problems},
J. Math. Anal. Appl. 525 (2023), no. 1, 127261.
\href{https://doi.org/10.1016/j.jmaa.2023.127261}{doi:10.1016/j.jmaa.2023.127261}.

\bibitem{BjorckGolub1973}
A. Bjorck and G. H. Golub,
\emph{Numerical methods for computing angles between linear subspaces},
Math. Comp. 27 (1973), no. 123, 579--594.
\href{https://doi.org/10.1090/S0025-5718-1973-0348991-3}{doi:10.1090/S0025-5718-1973-0348991-3}.

\bibitem{KoekoekLeskySwarttouw2010}
R. Koekoek, P. A. Lesky, and R. F. Swarttouw,
\emph{Hypergeometric Orthogonal Polynomials and Their $q$-Analogues},
Springer, Berlin, 2010.
\href{https://doi.org/10.1007/978-3-642-05014-5}{doi:10.1007/978-3-642-05014-5}.

\bibitem{KayalarWeinert1988}
S. Kayalar and H. L. Weinert,
\emph{Error bounds for the method of alternating projections},
Math. Control Signals Systems 1 (1988), no. 1, 43--59.
\href{https://doi.org/10.1007/BF02551235}{doi:10.1007/BF02551235}.

\bibitem{Yang2026}
J. Yang,
\emph{A Friedrichs angle between the Nyman--Beurling spaces and the Riemann hypothesis},
J. Math. Anal. Appl. 560 (2026), no. 1, 130494.
\href{https://doi.org/10.1016/j.jmaa.2026.130494}{doi:10.1016/j.jmaa.2026.130494}.
\end{thebibliography}
\end{document}
